% ===========================================================================
% ACT V -- Route Acquisition and Terrain Data  (Ch. 5)   Slides 41-48b
% Budget: 6 minutes. Supporting infrastructure, not the contribution -- keep
% it moving. Slides 48 and 48b were rewritten 2026-09-22 against branch
% surrogate-gns, where the corridor problem was diagnosed and largely fixed.
% ===========================================================================

\actdivider{5}{Routes and Terrain}

% --- 41. Why routes had to be built --------------------------------------
\begin{frame}{Rail Corridors Must Be Rebuilt from Map Data}
  \vspace{0.1em}
  \begin{columns}[c,onlytextwidth]
    \begin{column}{0.44\textwidth}
      \centering
      \includegraphics[width=\linewidth,height=0.80\textheight,keepaspectratio]%
        {../thesis_imgs_7_15/route_generator/ORM_base.png}
    \end{column}
    \begin{column}{0.53\textwidth}
      \begin{itemize}\setlength\itemsep{0.7em}
        \item Mainline corridors (orange) run through a dense tangle of yard and
              service track (brown), meeting at junctions.
        \item A corridor is not stored as one map object. It must be
              \textbf{assembled from many short segments}.
        \item Routes in this thesis come from \alert{real elevation and
              alignment data}, so grades and curves resemble an actual rail
              corridor rather than an invented profile.
      \end{itemize}
    \end{column}
  \end{columns}
\end{frame}

% --- 42. The pipeline -----------------------------------------------------
\begin{frame}{The Pipeline, in Five Stages}
  \vspace{1.2em}
  \centering
  \scalebox{0.74}{\routePipeline}

  \vspace{1.8em}
  \begin{minipage}{0.90\textwidth}\small
    Each stage is checked before the next is added, as with the simulator. Once
    an error in the grade reaches the simulator, it cannot be told apart from
    real terrain.
  \end{minipage}
\end{frame}

% --- 43. Acquisition ------------------------------------------------------
\begin{frame}{Stage 1: Corridor Extraction}
  \vspace{0.3em}
  \centering
  \includegraphics[width=0.90\linewidth,height=0.56\textheight,keepaspectratio]%
    {../thesis_imgs_7_15/route_generator/route_aq_page1.png}

  \vspace{0.5em}
  \begin{minipage}{0.92\textwidth}\footnotesize
    Map data is downloaded for an area first. Corridors are then recovered:
    yards and spurs are filtered out, the search expands outward up to a limit,
    fragments are joined, and corridors below a minimum length are discarded.
  \end{minipage}
\end{frame}

% --- 44. Resampling -------------------------------------------------------
\begin{frame}{Stage 1.5: Resampling to a Uniform Grid}
  \vspace{0.1em}
  \begin{columns}[c,onlytextwidth]
    \begin{column}{0.55\textwidth}
      \centering
      \includegraphics[width=\linewidth,height=0.70\textheight,keepaspectratio]%
        {../thesis_imgs_7_15/route_generator/route_resample.png}
    \end{column}
    \begin{column}{0.42\textwidth}
      \small
      A 49.1\,km corridor rewritten onto a uniform \textbf{10\,m grid}:
      4,920 points.

      \vspace{0.8em}
      \alert{713 of them are interpolated} across long source segments rather
      than sitting on a mapped vertex.

      \vspace{0.8em}
      These points stay flagged in later stages, so it is always clear where
      the geometry is estimated rather than mapped.
    \end{column}
  \end{columns}
\end{frame}

% --- 45. Elevation --------------------------------------------------------
\begin{frame}{Stage 2: Elevation, and Looking at It}
  \vspace{0.3em}
  \centering
  \includegraphics[width=0.78\linewidth,height=0.60\textheight,keepaspectratio]%
    {../thesis_imgs_7_15/route_generator/routes_elevated.png}

  \vspace{0.4em}
  \begin{minipage}{0.90\textwidth}\footnotesize
    The elevation is checked visually before grade is computed. A corridor
    joined incorrectly, or one that jumps between parallel tracks, is obvious in
    3D but \textbf{not visible in a grade plot}.
  \end{minipage}
\end{frame}

% --- 46. From elevation to grade -----------------------------------------
\begin{frame}{Stage 3: From Elevation to Grade}
  \vspace{0.3em}
  \centering
  \includegraphics[width=0.94\linewidth,height=0.46\textheight,keepaspectratio]%
    {../thesis_imgs_7_15/route_generator/route_elevation_profile.png}

  \vspace{0.6em}
  \begin{minipage}{0.92\textwidth}\small
    The line climbs from roughly 1430\,m to 1740\,m over 49\,km. Errors of about
    a metre in the elevation data are small compared with that climb. They are
    large compared with the \alert{elevation change between neighbouring 10\,m
    samples}, which is what grade measures. The pipeline therefore corrects
    bridges and tunnels first, then smooths and differentiates in one step.
  \end{minipage}
\end{frame}

% --- 47. The three fields -------------------------------------------------
\begin{frame}{The Three Fields the Simulator Consumes}
  \vspace{0.2em}
  \centering
  \scalebox{0.80}{\routeFieldsChart}

  \vspace{0.2em}
  \begin{minipage}{0.92\textwidth}\footnotesize
    Line~0, after the corridor fix on the next slide. The first 15\,km, through
    the urban area, is noisy. After that the grade stays near one percent.
  \end{minipage}
\end{frame}

% --- 48. The corridor set: diagnosis and fix -----------------------------
\begin{frame}{Finding and Removing a Bad Corridor}
  \vspace{0.1em}
  \begin{columns}[T,onlytextwidth]
    \begin{column}{0.54\textwidth}
      \vspace{0.6em}
      \rawElevQAChart
    \end{column}
    \begin{column}{0.43\textwidth}
      \vspace{0.2em}
      \small
      The first suspected cause was too little smoothing. \alert{That was not
      the cause.}

      \vspace{0.6em}
      Line~2's \emph{raw} elevation is wrong before any smoothing: 19.3\% of its
      10\,m steps jump more than a metre. Real track cannot climb a metre every
      10\,m over a fifth of its length, so \textbf{these samples are not on the
      track}.

      \vspace{0.6em}
      Smoothing cannot repair this. It only averages the incorrect values.
    \end{column}
  \end{columns}

  \vspace{0.4em}
  \begin{minipage}{0.96\textwidth}\footnotesize
    The fix is a quality check on the raw elevation, applied \emph{before}
    smoothing. Line~2 is rejected, the smoothing window increases from 200\,m to
    800\,m, and the grade limit drops from 4\% to \textbf{2.5\%}. The limit now
    works as an error alarm rather than a filter: 4\% is far steeper than freight
    main lines are built, so the old limit let a broken corridor pass for months.
  \end{minipage}
\end{frame}

% --- 48b. What the fix bought, and the wider corridor set ----------------
\begin{frame}{Results of the Fix, and a Larger Corridor Set}
  \vspace{0.2em}
  \begin{columns}[T,onlytextwidth]
    \begin{column}{0.46\textwidth}
      {\small\textbf{The Four Surviving Corridors}}
      \vspace{0.4em}

      {\footnotesize
      \begin{tabular}{@{}lrrr@{}}
        \toprule
        \textbf{Corridor} & \textbf{rms} & \textbf{p95} & \textbf{Max} \\
        \midrule
        Line 0 (49 km) & 0.86\% & 1.45\% & 2.50\% \\
        Line 3 (29 km) & 0.72\% & 1.25\% & 1.47\% \\
        Line 4 (27 km) & 0.39\% & 0.75\% & 1.23\% \\
        Line 5 (22 km) & 0.58\% & 0.87\% & 1.17\% \\
        \bottomrule
      \end{tabular}}

      \vspace{0.7em}
      {\footnotesize
      The worst overspeed falls from \textbf{67.8 to 25.7\,m/s}. The 245\,km/h
      freight trains came from Line~2's false grades.

      \vspace{0.4em}
      Runs over the limit by $>$5\,m/s are unchanged (13.3\% to 13.5\%).
      \alert{This has a second cause}: driving commands are generated
      without regard to the terrain.}
    \end{column}
    \begin{column}{0.50\textwidth}
      {\small\textbf{A Much Larger Corridor Set}}
      \vspace{0.5em}

      {\footnotesize The OpenStreetMap search failed five times in a row over
      4.5 hours. It was replaced by the USDOT/BTS North American Rail Network,
      which records how segments connect, labels yard track, and gives
      corridors real names (\emph{Moffat Tunnel, Union Pacific}) instead of
      \texttt{route\_line2}.}

      \vspace{0.6em}
      \begin{center}
        \stattile{95,936}{Main-Line Segments \textbullet\ 274,145 km}
      \end{center}

      \vspace{0.4em}
      {\footnotesize Next step: extract corridors from this network.}
    \end{column}
  \end{columns}
\end{frame}
